Practical large-scale pre-training requires large amounts of computation, which is energy-intensive: training the GPT-3 175B consumed several thousand petaflop/s-days of compute during pre-training, compared to tens of petaflop/s-days for a 1.5B parameter GPT-2 model (Figure \ref{figure:training flops}). This means we should be cognizant of the cost and efficiency of such models, as advocated by \cite{schwartz2019}. 
   
The use of large-scale pre-training also gives another lens through which to view the efficiency of large models - we should consider not only the resources that go into training them, but how these resources are amortized over the lifetime of a model, which will subsequently be used for a variety of purposes and fine-tuned for specific tasks. Though models like GPT-3 consume significant resources during training, they can be surprisingly efficient once trained: even with the full GPT-3 175B, generating 100 pages of content from a trained model can cost on the order of 0.4 kW-hr, or only a few cents in energy costs. Additionally, techniques like model distillation \cite{liu2019improving} can further bring down the cost of such models, letting us adopt a paradigm of training single, large-scale models, then creating more efficient versions of them for use in appropriate contexts. Algorithmic progress may also naturally further increase the efficiency of such models over time, similar to trends observed in image recognition and neural machine translation \cite{hernandez2020}.
